\chapter{FBI：Functional Body Interface——从具体身体到统一智能体身体语义}

\begin{chapteroverviewbox}
在前述 AgentOS 体系中，我们已经把一个完整智能体区分为认知内核与身体运行时两个彼此耦合但时间尺度不同的动力系统。认知内核处理目标、记忆、自我、推理与长期结构变化；身体运行时则直接面对现实世界中的传感器、执行器、设备状态、数字接口和高频局部动力学。问题由此变得十分具体：如果不同身体拥有完全不同的传感器、执行器、坐标系、控制协议和内部状态，一个持续存在的 Agent Kernel 如何在不随每一种身体重新设计自身的情况下获得统一的“身体经验”并发出统一意义上的行动意图？

本章提出 \textbf{FBI（Functional Body Interface，功能身体接口）} 作为这一问题的基础工程解。FBI 不是普通设备驱动层，也不是一个将传感器数据简单转换成文本的适配器，而是位于具体身体实现与 AgentOS 认知体系之间的\textbf{功能语义边界}。它一方面将身体特有的原始状态投影为统一的身体中心现实表示和功能影响表示，另一方面把认知层给出的抽象控制意图映射到该身体能够理解的执行参考空间，并向认知层显式声明该身体当前具有何种能力、约束、可信度与可达边界。

因此，本章建立的核心关系不是
\[
Sensor \rightarrow LLM,
\]
而是
\[
\boxed{
BodySpecificReality
\;\xleftrightarrow{\quad FBI\quad}\;
AgentInvariantBodySemantics
}
\]
其目标是在保持具体身体高效、局部、实时和可适应的同时，使上层 Agent Kernel 尽可能保持身体无关性。FBI 因而是 Body Scale 原理真正落入系统工程的第一道接口，也是后续 SGC、FCS、BPCC、KRA 以及多身体迁移能够成立的基础。
\end{chapteroverviewbox}
\ChapterArt{52}


\section{为什么智能体需要一个功能身体接口}

在没有身体抽象层的情况下，一个智能体与现实发生联系的最直接方式，是让认知系统直接接收摄像头图像、麦克风信号、电池电压、机器人关节角、浏览器 DOM、鼠标坐标、文件系统事件以及各种设备私有状态，并直接输出电机命令、鼠标轨迹、API 调用或系统调用。从短期演示角度看，这种方式并非不能工作；一个足够大的模型甚至能够在相当程度上学习各种输入输出格式之间的映射。然而，一旦我们的研究对象从一次性模型推理转向持续存在、可以更换身体、长期成长的 Agent，这种设计立即暴露出根本缺陷：认知结构被大量偶然的身体实现细节污染了。

设智能体可以使用的身体集合为
\[
\mathcal B
=
\{B_1,B_2,\ldots,B_n\}.
\]
每一种身体 \(B_i\) 都具有自己的原始观测空间
\[
\mathcal Y_i^{raw},
\]
身体内部状态空间
\[
\mathcal X_i^{body},
\]
以及原始执行空间
\[
\mathcal U_i^{raw}.
\]
如果 Agent Kernel 直接面对这些空间，那么对于不同身体，原则上需要学习不同的闭环：
\[
\pi_i:
\mathcal Y_i^{raw}
\times
\mathcal X_i^{body}
\rightarrow
\mathcal U_i^{raw}.
\]
此时，即使两个身体完成的是语义上完全相同的行为，例如“走到桌子旁”“打开文件”“把对象交给另一个主体”，它们在内核看来仍然表现为截然不同的输入输出问题。这样形成的不是一个具有多个身体的持续智能体，而是若干被身体接口切割开的局部智能实现。

我们希望建立的关系恰好相反。认知内核应该更多处理诸如“对象在哪里”“身体是否可以到达”“当前身体承受何种约束”“动作是否已经完成”“当前还有多少控制余量”等\textbf{功能意义}，而不是长期记忆某一类摄像头的像素排列、某一机器人关节的编号或者某个浏览器版本的事件格式。因此我们引入一个相对稳定的标准语义空间
\[
\mathcal Z_B,
\]
并要求每一种具体身体存在一个映射
\[
\Phi_i:
(\mathcal Y_i^{raw},\mathcal X_i^{body})
\rightarrow
\mathcal Z_B,
\]
同时存在与执行方向相对应的映射
\[
\Psi_i:
\mathcal U_B
\rightarrow
\mathcal U_i^{body},
\]
其中 \(\mathcal U_B\) 是 AgentOS 统一的身体控制参考空间，而不是某个具体执行器的电压或扭矩空间。

于是，真正希望训练和保持稳定的是较高层关系
\[
\Pi:
\mathcal Z_B
\rightarrow
\mathcal U_B,
\]
而不同身体之间的差异尽可能被约束在
\[
(\Phi_i,\Psi_i)
\]
以及相应的局部身体动力学之中。由此形成：
\[
\boxed{
\pi_i
\approx
\Psi_i
\circ
\Pi
\circ
\Phi_i
}
\]
这一分解是 FBI 的第一性意义。

这里必须特别强调，“统一接口”并不是要求不同身体拥有相同的感官或相同的动作集合。一个桌面数字 Agent 没有机械臂，一个轮式机器人不能像人类一样跨越楼梯，一个机械臂也没有浏览器中的“点击链接”动作。FBI 所做的不是伪造这种同一性，而是把\textbf{能够统一的语义统一起来，把不能统一的能力差异显式暴露出来}。因此身体抽象不是掩盖身体差异，而是将差异从混乱的实现细节提升为可以被 Agent 推理的能力结构。

我们由此获得一条后续贯穿 Body Runtime 的原则：
\[
\boxed{
BodyAbstraction
\neq
BodyErasure
}
\]
FBI 消除的是无意义的实现耦合，而不是身体本身。恰恰相反，当底层差异被正确抽象后，身体的真实能力边界才第一次能够成为上层认知的显式对象。


\section{FBI 的系统位置与功能边界}

在整个 AgentOS 中，我们将完整智能体写为
\[
\boxed{
AgentSystem
=
AgentKernel
\bowtie
BodyRuntime
}
\]
其中符号 \(\bowtie\) 强调二者不是普通单向 API 调用关系，而是持续闭环耦合。Agent Kernel 的主要状态包含工作记忆 \(h\)、情景经验 \(E\)、长期结构 \(\Theta\)、自我结构 \(\Psi\)、生命周期慢变量以及相关认知调度状态；Body Runtime 则维护与当前身体和现实世界直接相关的快速状态。

在这一结构中，可以将 Body Runtime 抽象表示为
\[
\boxed{
BodyRuntime
=
FBI
\bowtie
LocalDynamics
\mid
SafetyBoundary
}
\]
其中 \(LocalDynamics\) 后续可以由高频预测与控制组件承担，而 FBI 负责公共身体语义和接口约束。这里必须明确：\textbf{FBI 本身不是完整低层控制器。}如果把传感器解码、语义融合、局部预测、运动规划、伺服控制、身体适应和认知接口全部塞进 FBI，那么 FBI 很快会变成一个无法稳定定义的巨型模块，也失去作为长期 ABI 的意义。

因此 FBI 应承担的是一组更稳定的边界职责：
\[
\fitdisplay{
FBI
=
\{
Acquisition,
Normalization,
Grounding,
StatePublication,
ReferenceAcceptance,
CapabilityPublication
\}.
}
\]
其中 Acquisition 表示从设备、操作系统或数字环境获得身体数据；Normalization 表示单位、时间戳、坐标和质量信息的规范化；Grounding 表示把设备特有状态绑定到 Agent 可理解的身体与现实语义；StatePublication 表示向上层发布统一身体状态；ReferenceAcceptance 表示接受上层的控制参考，而不是未经约束的任意底层命令；CapabilityPublication 则负责显式说明该身体能做什么、不能做什么以及当前能力是否退化。

从时间尺度上看，这种边界尤其重要。设底层伺服控制周期为
\[
\tau_{servo},
\]
身体局部预测控制周期为
\[
\tau_{body},
\]
FBI 公共语义更新周期为
\[
\tau_{FBI},
\]
认知工作循环为
\[
\tau_h.
\]
通常应满足
\[
\boxed{
\tau_{servo}
\leq
\tau_{body}
<
\tau_{FBI}
\lesssim
\tau_h
}
\]
但它们并不要求严格固定比例。一个高速碰撞传感器可能以毫秒级运行，而认知层只需接收“当前碰撞风险快速上升”这一具有稳定意义的状态。FBI 的作用之一，正是在多个物理时间尺度之间建立可供认知使用的语义更新边界。

我们还应区分三个不同概念。第一是 \textbf{Body Driver}，它直接面对具体硬件或数字设备协议；第二是 \textbf{FBI}，它建立稳定身体语义；第三是 \textbf{Agent Kernel}，它依据这些语义形成目标、推理和决策。三者可以写为
\[
RawDevice
\rightarrow
BodyDriver
\rightarrow
FBI
\rightarrow
AgentKernel.
\]
反向则为
\[
AgentKernel
\rightarrow
FBI
\rightarrow
BodyRuntime
\rightarrow
BodyDriver
\rightarrow
Actuator.
\]

这一结构类似计算机操作系统中设备驱动与抽象设备接口之间的区别，但又比传统 HAL 更进一步。传统硬件抽象层通常只要求“设备以统一调用方式工作”；FBI 还要求不同身体把状态转换成\textbf{智能体能够长期学习和迁移的功能语义}。因此 FBI 是硬件抽象、身体语义抽象和认知 ABI 三者交叉形成的新层次。


\section{Body Driver：保留身体差异的最底层适配}

FBI 若要保持长期稳定，就必须避免直接吸收每一种硬件的偶然实现。因此，具体身体差异首先由 Body Driver 层处理。Body Driver 与传统设备驱动存在部分同构，但其服务对象不再只是操作系统，而是一个持续存在的智能体身体运行时。

对某一身体 \(B_i\)，我们可以把原始接口写成
\[
D_i
=
(
Y_i,
X_i,
U_i,
\Sigma_i
),
\]
其中 \(Y_i\) 表示外部传感器输出，\(X_i\) 表示设备内部状态，\(U_i\) 表示可执行命令，\(\Sigma_i\) 表示设备协议、时序和异常状态集合。Body Driver 首先建立
\[
\mathcal D_i:
D_i
\rightarrow
D_i^{norm},
\]
把私有协议转换成规范化设备对象。

这种规范化至少应覆盖四个方面。第一是\textbf{单位规范化}。速度、距离、温度、电量、关节角以及计算资源等不能让上层长期依赖设备私有编码。第二是\textbf{时间规范化}。任何感知状态都必须保留时间戳、采样周期与延迟估计，否则来自不同设备的状态在融合时会形成虚假的同时性。第三是\textbf{坐标规范化}。物理机器人存在世界坐标、身体坐标、传感器坐标与执行器坐标；数字身体同样存在屏幕坐标、DOM 层级、文件路径和应用上下文。第四是\textbf{质量规范化}，即状态不仅包含数值，还应带有置信度、有效期、异常标志与来源信息。

因此，一个规范化感知项不应只是
\[
y_k=v_k,
\]
而更接近
\[
\hat y_k
=
(
v_k,
t_k,
frame_k,
unit_k,
q_k,
source_k
),
\]
其中 \(q_k\) 表示质量或置信属性。

Body Driver 的第二项重要职责是保持设备能力和故障的真实性。例如，如果某个机械臂第三关节失效，上层不应继续看到一个虚假的“完整机械臂”；如果浏览器页面权限被撤销，数字 Agent 也不应继续认为自己可以执行全部交互。因此设备状态必须向 FBI 暴露 capability degradation：
\[
C_i(t)
=
C_i^{nominal}
-
\Delta C_i(t).
\]
这使身体能力从一个静态配置文件变成一个随现实动态变化的状态。

我们尤其需要避免一个常见设计错误：让 Body Driver 直接产生复杂认知解释。例如摄像头驱动不应该输出“这是一个危险的人”，鼠标驱动也不应该决定“这个按钮值得点击”。Driver 层应该尽可能保持确定、局部、可验证。更高层语义由 FBI 的语义融合部分和 Body Runtime 其他组件完成。这样一旦高层模型变化，底层设备适配仍能保持稳定。

这一原则可以写成：
\[
\boxed{
DriverKnowsDevice,
FBIKnowsBodyMeaning,
KernelKnowsTaskMeaning.
}
\]

这种分层在长期智能系统中特别重要，因为身体通常比认知模型更容易发生替换。我们可能升级摄像头、改变机械臂型号、切换浏览器、迁移到新的操作系统。如果每一次硬件变化都迫使长期记忆 \(\Theta\) 与自我结构 \(\Psi\) 大规模重新学习，那么所谓持续智能就无法真正成立。Body Driver 因而承担着保护高层认知连续性的第一道工程职责。


\section{Body State Core：身体并不等于工作记忆}

在早期 Agent 架构中，一个常见隐含假设是：只要把环境和设备状态放入大模型上下文，模型就“拥有了身体”。这种方式混淆了身体状态与认知工作状态。为了避免这一问题，FBI 必须维护独立的 Body State Core。

我们定义身体状态
\[
X_B(t)
\]
与认知工作状态
\[
h(t)
\]
满足
\[
\boxed{
X_B(t)\neq h(t).
}
\]
身体状态可以非常庞大，并且许多变量以远高于认知循环的频率变化。例如机器人可能拥有数百个关节、电流、温度、触觉和局部动态变量；桌面 Agent 可能同时面对数千个 UI 对象、网络连接、进程与文件状态。让这些状态全部持续进入 \(h(t)\)，将直接违背前述工作记忆有限容量原则。

因此 Body State Core 应保存比工作记忆更完整的局部身体状态：
\[
X_B(t)
=
\{
X^{kin},
X^{dyn},
X^{device},
X^{resource},
X^{health},
X^{environment-local}
\}.
\]
它是身体运行时自己的状态空间，而 FBI 根据当前认知需求、变化显著性和接口契约，从中形成面向 Agent Kernel 的公共投影：
\[
Z_B(t)
=
\Phi_{pub}
\bigl(
X_B(t)
\bigr).
\]
这里的 \(Z_B(t)\) 并不是简单摘要，而是一组具有持续语义的身体公共状态。

我们可以进一步把身体状态分为三个层级：
\[
X_B
=
X_B^{raw}
\oplus
X_B^{derived}
\oplus
X_B^{public}.
\]
其中 \(X_B^{raw}\) 是传感器和设备原始状态；\(X_B^{derived}\) 是 Body Runtime 内部估计出的速度、趋势、稳定性、接触状态、资源余量等；\(X_B^{public}\) 是经过 FBI 暴露给 Agent Kernel 的稳定语义状态。这个分层保证了上层可以获得足够信息，却不必承担底层每一个变量的计算压力。

更重要的是，Body State Core 应具有\textbf{持续性}。如果认知层暂时没有关注某个身体变量，该变量仍然继续存在和演化。例如 Agent 没有把“电池温度”放入工作记忆，并不意味着身体温度停止变化。这正对应我们此前形成的原则：
\[
\boxed{
MostAgentDynamics\notin h.
}
\]
成熟智能的一个重要特征，就是大量正常动力学能够在局部闭环中独立运行，而只有与当前目标、异常、风险或显著变化相关的结构才被提升到认知层。

从控制论角度，可以把 Body State Core 看成一个持续估计的局部状态：
\[
\dot X_B
=
F_B(X_B,u_B,w),
\]
其中 \(u_B\) 是身体控制输入，\(w\) 表示现实扰动。FBI 的公共状态只是其投影：
\[
Z_B
=
H_B(X_B).
\]
这解释了为什么认知层不能把自己看到的身体摘要误认为身体真实状态本身。任何公共语义都具有采样、压缩和延迟，因此还应带有认识论属性，例如：
\[
Z_B^k
=
(
value,
confidence,
age,
source
).
\]

这种设计实际上为 AgentOS 保持了非常重要的认识论谦逊：Agent 对自己的身体也不是“全知”的，它同样通过一个有限的接口不断估计和更新。身体状态不是一个透明数据库，而是一个必须持续通过现实反馈校准的动力对象。


\section{SGC Interface：把身体所处现实变成统一可行动语义}

FBI 向 Agent Kernel 提供的第一类核心公共状态，是关于“外部现实及身体在现实中的位置”的结构表示。在本书中，这一接口统一称为 SGC（Semantic Graph Coordinate Space）。SGC 的完整结构将在后续章节专门讨论，本节只从 FBI 的接口职责出发说明它为什么必须存在，以及 FBI 应当向这一空间提供什么。

设某一时刻身体获得原始外部观测
\[
Y^{ext}(t).
\]
如果直接把 \(Y^{ext}\) 输入认知内核，那么不同身体仍然表现为不同模态：摄像头是像素，激光雷达是点云，浏览器是 DOM，文件系统是节点树，Shell 是字符流。认知层真正需要的却往往是跨身体可复用的结构：
\[
Entity,
Relation,
Position,
State,
Affordance,
Uncertainty.
\]
因此 FBI 需要建立：
\[
\Phi_{SGC}:
(Y^{ext},X_B)
\rightarrow
Z_{SGC}.
\]

这里的关键不是把所有东西“翻译成文字”，而是建立一个身体中心的结构世界。一个对象可以表示为
\[
o_i
=
(
id_i,
g_i,
a_i,
s_i,
r_i,
e_i
),
\]
其中 \(g_i\) 是几何或位置状态，\(a_i\) 是相对稳定属性，\(s_i\) 是动态状态，\(r_i\) 是与其他结构的关系，\(e_i\) 是认识论元数据。于是不同身体即使感知方式不同，也可以在更高层形成类似的对象结构。

例如物理机器人通过视觉和深度相机识别出一扇门，数字 Agent 通过网页结构识别出一个“确认按钮”。二者在底层完全不同，但从智能体功能角度，都可能具有：
\[
Affordance:
CurrentState
\rightarrow
PossibleAction.
\]
FBI 的任务不是声称“门”和“按钮”是同一种物体，而是让上层认知能够使用共同的结构关系：它们都是当前身体可以作用于现实的可行动对象。

SGC 中另一个关键原则是
\[
\boxed{
Body\in SGC.
}
\]
也就是说，Agent 不能仅获得一个第三人称世界模型。世界必须包含“我在哪里”“我的作用范围在哪里”“哪些对象相对于我可达”等身体相关关系。如果没有这一点，所谓 world model 仍然只是脱离主体的地图，而不是可供行动的现实。

因此 FBI 至少需要向 SGC 提供三个类别的数据。第一类是外部实体和场的规范化状态；第二类是身体自身的几何与行动位置；第三类是身体与环境之间的关系，例如可达、接触、遮挡、占用、拥有、控制、权限等。

值得注意的是，SGC 并不要求 FBI 完成全部高级语义理解。FBI 可以负责将底层身体状态稳定落入标准语义槽位，而更复杂的对象身份、社会意义、因果关系可以由更高层模型进一步补充。这样形成的是一个逐层丰富的结构，而不是要求接口层一次解决整个世界模型问题。

因此从 FBI 角度看，SGC 最重要的意义是：
\[
\boxed{
BodySpecificPerception
\rightarrow
AgentStableActionableReality.
}
\]
它使一个 Agent 换身体以后，不必重新学习“世界存在对象、关系、位置、状态和可行动性”这些高层结构，只需要重新学习新的身体如何把现实投影进这一公共空间。


\section{FCS Interface：把身体受到的作用变成统一功能状态}

如果 SGC 主要回答“现实中有什么”，那么仅有 SGC 仍然不足以构成完整身体接口。智能体还必须知道现实正在如何作用于自己的身体。电量下降、温度升高、机械结构受力、视觉输入过载、网络延迟增加、控制不稳定、资源耗尽或某个数字系统权限即将失效，这些状态未必首先表现为一个外部对象，却会直接改变智能体当前行动的可行性和优先级。

因此 FBI 的第二类核心公共接口是 FCS。其完整定义将在后续章节展开；本节只讨论它作为 FBI 输出的功能意义。我们可以把 FCS 抽象为：
\[
\Phi_{FCS}:
(X_B,Y^{int},Y^{ext})
\rightarrow
Z_{FCS},
\]
其中 \(Z_{FCS}\) 表示现实与身体当前产生的功能性影响场。

必须特别注意，FCS 中的 ``Feeling'' 是工程意义上的 functional influence，而不是对主观感受的断言。一个机器人可以具有
\[
ThermalPressure=0.8
\]
并不意味着它在现象意识意义上“感觉热”；它仅意味着当前热状态已经对行动、安全或资源调度构成强约束。类似地，一个数字 Agent 可以具有
\[
ComputePressure,
LatencyPressure,
PermissionRisk,
\]
而这些变量同样可以进入统一的功能影响结构。

一个典型 FCS 通道可以表示为
\[
f_k(t)
=
(
I_k,
\dot I_k,
q_k,
m_k
),
\]
其中 \(I_k\) 是当前影响强度，\(\dot I_k\) 表示变化趋势，\(q_k\) 表示可信度，\(m_k\) 表示有效性或空间掩码。这样的设计比单一告警信号丰富，因为控制系统通常不仅需要知道“当前高不高”，还需要知道“正在快速变坏还是正在恢复”。

FBI 在这里承担的关键任务，是把各种身体私有内部状态映射为相对稳定的功能通道。例如不同机器人使用不同电池技术，但对于 Agent Kernel 来说，更持久的变量可能是：
\[
EnergyReserve,
EnergyPressure,
ExpectedOperatingTime.
\]
不同数字身体的 CPU、GPU、网络和 API quota 也可以在一定程度上映射为资源压力，而不是让认知层记忆每种硬件的所有底层指标。

这一抽象仍然不能过度。不同身体可能具有真正独特的功能影响，FCS 因而应该具有固定核心通道和可扩展专用通道，而不是强迫全部身体压进极少数变量。我们可以写：
\[
Z_{FCS}
=
Z_{core}
\oplus
Z_{body-specific}.
\]
核心部分保证迁移，专用部分保留身体个性。

从更高层看，FCS 与 SGC 构成 FBI 的两种互补投影：
\[
\boxed{
SGC:\ WhatIsThere?
}
\]
\[
\boxed{
FCS:\ HowIsRealityAffectingThisBody?
}
\]
前者主要建立外部可行动现实，后者建立身体约束与功能状态。它们共同使 Agent Kernel 不必直接处理全部原始传感器，却仍然能够在长期学习中形成关于“世界”和“我在世界中受到什么作用”的稳定结构。

这种二元设计对于持续身份尤其重要。Agent 可以更换摄像头、机械臂甚至整个载体，但只要新的身体仍然能够重新建立 SGC/FCS 公共语义，上层长期认知结构就存在继续工作的可能。由此，FBI 开始真正把前述 Body Scale 原理从抽象原则转化为可以工程测量的接口迁移能力。


\section{Control Reference：认知意图不应直接等于执行器命令}

FBI 不仅需要解决身体向认知的输入，还必须处理认知向身体的输出。如果 Agent Kernel 直接产生执行器级命令，例如每个轮子的速度、机械臂各关节扭矩、鼠标每毫秒的像素位移，那么高层认知就会重新与具体身体绑定，并承担其不应承担的高频控制负担。因此 FBI 的输出方向必须建立一种介于“抽象 Goal”和“底层 actuator command”之间的稳定控制语言，即 Control Reference。

我们定义高层控制参考为
\[
r_t
\in
\mathcal R_B.
\]
它不是要求身体立即执行一个确定微动作，而是描述认知层希望身体进入的可行目标区域。一个控制参考可以形式化为：
\[
r_t
=
(
g_t,
\mathcal E_t,
\mathcal C_t,
p_t,
T_t
),
\]
其中 \(g_t\) 是目标状态或目标关系，\(\mathcal E_t\) 是允许的误差包络，\(\mathcal C_t\) 是约束集合，\(p_t\) 是优先级或行为模式，\(T_t\) 是相关时间窗口。

例如“移动到桌子旁边”不需要 Agent Kernel 指定每一个关节轨迹，而可以给出：
\[
Goal:
distance(body,table)<d^\ast,
\]
以及：
\[
CollisionRisk<\epsilon,
\qquad
EnergyPressure<\eta.
\]
身体运行时再根据当前身体结构和现实条件完成局部路径与控制。这正体现了：
\[
\boxed{
HigherLevelSetsFeasibleRegion,
\quad
LowerLevelOptimizesInsideIt.
}
\]

对于数字身体也是一样。高层认知可以表达“打开目标文档并进入编辑状态”，而不必长期绑定到具体操作系统中的鼠标坐标序列。如果当前应用支持语义 API，Body Runtime 可以调用 API；如果只能通过 GUI，则可以使用视觉和鼠标控制。只要最终满足相同的控制参考，上层 Agent-CSO 就可以保持较高程度连续。

FBI 在接收 Control Reference 时必须进行第一层契约检查。设身体当前能力集合为
\[
\mathcal C_B(t),
\]
控制请求需要的能力为
\[
\mathcal C_r.
\]
只有满足
\[
\mathcal C_r
\subseteq
\mathcal C_B(t)
\]
时，请求才具有直接执行可能性。否则 FBI 不应静默失败，而应返回结构化不可达结果：
\[
Failure
=
(
MissingCapability,
ConstraintViolation,
Uncertainty,
Alternative
).
\]
这使“身体做不到”成为可以被认知层推理的明确事实，而不是执行失败以后才出现的模糊错误。

此外，Control Reference 必须与安全权限分离。认知层可以表达目标，但不能自动获得绕过身体安全边界的能力。即使 Kernel 希望身体进入某个状态，Body Runtime 仍然必须在现实、能力和安全约束内执行。因此：
\[
\boxed{
Intent
\neq
Authority
\neq
Execution.
}
\]

这一点对于长期 Agent 极为重要。如果认知规划直接拥有所有执行器的无限写权限，那么任何短期推理错误都有可能立即变成不可逆现实后果。FBI 因而既是语义接口，也是第一层执行契约边界。

从更一般的角度说，Control Reference 把“我要做什么”与“这个身体如何做到”分离。前者属于长期可以迁移的认知结构，后者属于身体特有的局部技能。只有建立这种分离，Agent 才可能真正经历：
\[
\boxed{
SameAgent,
DifferentBody.
}
\]


\section{Body Capability Descriptor：让身体能力成为认知可以推理的对象}

如果 FBI 只统一输入和输出，却不显式描述身体能力，Agent Kernel 仍然会产生一个严重问题：它不知道当前身体究竟能做什么。传统软件往往把设备能力藏在 API 文档或驱动实现中，而持续 Agent 需要把这些能力本身变成可观测、可比较、可更新的认知结构。因此 FBI 必须提供 Body Capability Descriptor。

可以将某个身体的能力描述写为
\[
\mathcal K_B(t)
=
(
\mathcal P,
\mathcal A,
\mathcal R,
\mathcal L,
\mathcal S,
\mathcal Q
),
\]
其中 \(\mathcal P\) 表示感知能力，\(\mathcal A\) 表示行动能力，\(\mathcal R\) 表示资源条件，\(\mathcal L\) 表示限制，\(\mathcal S\) 表示安全与权限边界，\(\mathcal Q\) 表示当前能力质量。

每项能力又不应只是 Boolean：
\[
can\_move=true.
\]
更实际的形式是：
\[
c_k
=
(
domain,
range,
precision,
latency,
reliability,
cost,
authority
).
\]
例如一个机械臂的抓取能力必须包含工作空间、最大负载、定位精度、预计延迟、失败率和当前安全等级；一个数字 Agent 的“写文件”能力则可能包含允许目录、权限、文件大小、应用上下文和审计要求。

于是，身体能力本身构成一个动态可行域：
\[
\mathcal A_B(t)
=
\left\{
a\;|\;
g_j(a,X_B(t))\leq 0,
\ j=1,\ldots,m
\right\}.
\]
Agent Kernel 的规划不应该在无限动作空间中搜索，而应在 FBI 暴露的当前可行域之上形成策略：
\[
\pi^\ast
=
\arg\max_{\pi\in\Pi(\mathcal A_B)}
J(\pi).
\]
这样，身体能力从“执行以后才知道”变成“决策之前就可以参与推理”。

Capability Descriptor 还承担身体迁移中的关键角色。设 Agent 从身体 \(B_i\) 切换到 \(B_j\)，我们并不要求：
\[
\mathcal K_{B_i}
=
\mathcal K_{B_j}.
\]
真正需要计算的是能力映射：
\[
M_{ij}:
\mathcal K_{B_i}
\rightarrow
\mathcal K_{B_j}.
\]
其中部分能力可以直接迁移，部分需要重新 grounding，部分根本不存在。于是换身体可以被明确分解成：
\[
SharedCapability,
RegroundedCapability,
MissingCapability,
NewCapability.
\]

这种结构与我们此前提出的 Body Scale 直接相关。Body Scale 越高，并不是两个身体越相似，而是：
\begin{statementbox}
已有高层认知结构在新的 Capability Descriptor 下可复用的比例越高。
\end{statementbox}
因此未来可以定义概念性的身体迁移率：
\[
T_{B_i\rightarrow B_j}
=
\frac{
\text{可保持高层语义的能力结构}
}{
\text{目标身体所需能力结构}
}.
\]

Capability Descriptor 也使 Agent 能够形成关于自身身体的显式认识。例如一个 Agent 可以知道“当前身体无法上下楼梯”“视觉传感器正在退化”“浏览器没有文件系统写权限”。这些事实可以进入工作记忆 \(h\)，成为真正的 Self-in-World 结构，而不需要认知模型通过连续失败间接猜测自己的能力。

从这一角度看，FBI 不只是“让身体可用”，它还第一次把身体从黑箱工具提升成 Agent 可以理解和推理的\textbf{自身能力结构}。


\section{FBI 的统一数学模型与接口不变量}

经过前面的分解，我们现在可以给 FBI 一个更严格的形式定义。对具体身体 \(B_i\)，定义：
\[
B_i
=
(
\mathcal X_i,
\mathcal Y_i,
\mathcal U_i,
\mathcal K_i
),
\]
分别代表身体状态空间、观测空间、执行空间和能力空间。FBI 为该身体建立一组映射：
\[
FBI_i
=
(
\Phi_i^{state},
\Phi_i^{SGC},
\Phi_i^{FCS},
\Psi_i^{ctrl},
\Gamma_i^{cap}
).
\]

其中
\[
\Phi_i^{state}:
(\mathcal X_i,\mathcal Y_i)
\rightarrow
\mathcal X_B^{pub}
\]
产生公共身体状态；

\[
\Phi_i^{SGC}:
(\mathcal X_i,\mathcal Y_i)
\rightarrow
\mathcal Z_{SGC}
\]
产生行动相关现实结构；

\[
\Phi_i^{FCS}:
(\mathcal X_i,\mathcal Y_i)
\rightarrow
\mathcal Z_{FCS}
\]
产生身体功能影响结构；

\[
\Psi_i^{ctrl}:
\mathcal R_B
\rightarrow
\mathcal R_i^{local}
\]
把统一 Control Reference 转换到本地控制问题；

\[
\Gamma_i^{cap}:
\mathcal X_i
\rightarrow
\mathcal K_B(t)
\]
持续发布当前身体能力。

于是 FBI 的公共 ABI 可以概括为：
\[
\boxed{
FBI^{pub}
=
(
SGC,
FCS,
BodyState,
Capability,
ControlReference,
Status
).
}
\]

真正重要的是，这个 ABI 必须满足一组接口不变量。第一是\textbf{语义不变量}：更换身体以后，同一公共字段的功能意义不能任意改变。第二是\textbf{时间不变量}：所有公共状态必须具备明确的时间语义、延迟和有效期。第三是\textbf{来源不变量}：观测状态、推导状态与预测状态不能无标记混合。第四是\textbf{控制不变量}：Control Reference 表示目标和约束，而不把某个具体执行器的私有动作泄漏到高层。第五是\textbf{能力真实性不变量}：FBI 不得声明身体当前实际上不具有的能力。第六是\textbf{现实再进入不变量}：执行结果必须通过现实观测重新进入 FBI，而不能由“命令已发出”直接推断“动作已成功”。

特别是最后一点，可以写为：
\[
\boxed{
ActionIssued
\not\Rightarrow
ActionSucceeded.
}
\]
若发出控制参考 \(r_t\)，真正的状态更新必须经过：
\[
r_t
\rightarrow
BodyDynamics
\rightarrow
Reality
\rightarrow
Observation
\rightarrow
FBI
\rightarrow
Z_{t+1}.
\]
这条闭环保证 AgentOS 中最终权威属于现实，而不是属于内部预测模型。

为了评价 FBI 是否实现成功，我们还可以定义几个工程指标。首先是语义保持率：
\[
R_{sem},
\]
衡量不同身体接入后高层结构是否保持相同意义。其次是身体切换重学习成本：
\[
C_{relearn}(B_i\rightarrow B_j).
\]
第三是接口延迟：
\[
L_{FBI}.
\]
第四是状态完整性：
\[
I_{state}.
\]
第五是能力校准误差：
\[
E_{cap}
=
d(
\mathcal K_B^{declared},
\mathcal K_B^{real}
).
\]
理想情况下，我们希望：
\[
R_{sem}\uparrow,
\qquad
C_{relearn}\downarrow,
\qquad
L_{FBI}\downarrow,
\qquad
E_{cap}\downarrow.
\]

于是 Body Scale 终于获得可以工程化讨论的基础：一个好的身体接口，不是把所有身体变成同一个身体，而是最大程度保持智能体高层认知结构与具体身体实现之间的可分离性。

我们可以把本章的最终原则压缩为：
\[
\boxed{
\frac{\partial Architecture_{Agent}}
{\partial Body_i}
\rightarrow 0
}
\]
这里当然不是严格数学意义上要求导数等于零，而是表示一种架构设计目标：\textbf{随着身体实现发生变化，Agent 的长期认知核心应该尽可能少发生结构性重写。}

因此，FBI 的最终定义可以写成：

\begin{definition*}
FBI 是具体身体动力学与持续智能主体之间的稳定功能语义边界。它把身体特有的感知、内部状态、执行能力和约束投影为 Agent 可长期理解的公共身体结构，同时把高层控制意图限制性地映射回具体身体，使“同一个 Agent 使用不同身体”成为可以工程实现而不仅是哲学假设的问题。
\end{definition*}

这也是为什么在 AgentOS 的工程结构中，FBI 必须先于更复杂的身体预测、技能编译和跨身体适应机制被稳定定义：如果身体与认知之间不存在一个长期稳定的公共语义契约，那么任何后续的身体学习都会不断反向污染认知内核，所谓 Body Scale 与持续身份也就失去了真正的工程基础。
